\begin{frame}[fragile]{Topologische Quantisierung der Perihel-Präzession (Gate $G_5$)}
    \begin{columns}[T]
        \column{0.33\textwidth}
        \textbf{1. Geodätische Baseline}
        \vspace{0.2cm}
        \begin{itemize}
            \small
            \item Hyperbolische Modulfläche $\mathcal{X}_{\text{mod}} = PSL(2,\mathbb{Z})\backslash\mathbb{H}^2$.
            \item Fermat-Primzahl $p = 5 \equiv 1 \pmod 4$ in $\mathbb{Z}[i]$.
            \item Ungestörte Bahnlänge:
            $$\ell(5) = \ln(5) \approx \mathbf{1{,}609437}$$
        \end{itemize}

        \column{0.33\textwidth}
        \textbf{2. Weyl-Dämpfung}
        \vspace{0.2cm}
        \begin{itemize}
            \small
            \item Drift als hochenergetische Weyl-Mode im Co-Sektor.
            \item EABC-Grenzrate:
            $$\lambda_{\text{Weyl}} = \frac{3}{16} = 0{,}1875$$
            \item Dämpfung nach einem Umlauf:
            $$\frac{A(\ell(5))}{A(0)} = 5^{-3/16} \approx \mathbf{0{,}7395}$$
        \end{itemize}

        \column{0.33\textwidth}
        \textbf{3. Ptolemäisches Locking}
        \vspace{0.2cm}
        \begin{itemize}
            \small
            \item Fermat-Gauss-Phasen auf $U(1)$ ($k=1..4$).
            \item Sehnenviereck nach Satz von Ptolemäus:
            $$d_{13}d_{24} \equiv d_{12}d_{34} + d_{14}d_{23}$$
            \item Numerische Identität:
            $$\mathbf{0{,}8000 \equiv 0{,}8000}$$
        \end{itemize}
    \end{columns}

    \vspace{0.4cm}
    \centering
    \fbox{\parbox{0.95\textwidth}{\centering\footnotesize
    \textbf{Physikalisches Fazit:} Das Ptolemäische Korsett ist unelastisch. Jede kontinuierliche Verzerrung verletzt die algebraische Invarianz $\rightarrow$ Die Perihel-Präzession wird topologisch auf diskrete Weyl-Sektoren quantisiert.
    }}
\end{frame}